
\section{Literature-audit methodology}
\subsection{Queries and retrieval}
197 PubMed records were individually fetched through NCBI E-utilities
(esearch, relevance-ranked, then esummary per record) with these queries:
\begin{itemize}
\item \texttt{malaria parasite detection deep learning blood smear}
\item \texttt{automated malaria microscopy diagnosis}
\item \texttt{Plasmodium image classification machine learning}
\item \texttt{malaria thin smear image analysis computer aided}
\item \texttt{malaria red blood cell segmentation detection}
\end{itemize}
Retrieval was capped at 150 records per query; deduplication by PMID; the
shared suite audit's malaria subset was unioned in and deduplicated again.
Every record carries its PMID, title, journal, and date in
\texttt{results/lit\_audit\_malaria.json} --- accession-level evidence, not
a count claim.
\subsection{Screening and clustering}
Records were screened by title and assigned to thematic clusters by
keyword analysis (deep learning, CAD/segmentation, classical ML,
dataset/benchmark, clinical context, label noise/robustness). The
screen's limits are stated in the survey: title-keyword clustering can
misfile borderline records, PubMed coverage is not the whole literature,
and the audit date bounds recency claims. No record was excluded from the
ledger for being inconvenient to our novelty claim --- the label-noise
cluster's tiny size \emph{is} the finding.
